\documentclass[fleqn]{article}

\usepackage{aima-slides}
\usepackage[utf8]{inputenc}
\usepackage[T5]{fontenc}
\usepackage{lmodern}

\begin{document}

\begin{huge}
\titleslide{Giải quyết vấn đề và tìm kiếm}{Chương 3, Phần 1--5}

%%%%%%%%%%%% Slide %%%%%%%%%%%%%%%%%%%%%%%%%%%%%%%%%%%%%%%%%%%%%%%%%%%
\heading{Nội dung}

\blob Tác nhân giải quyết vấn đề

\blob Các loại bài toán

\blob Khởi tạo bài toán

\blob Các bài toán ví dụ

\blob Các thuật toán tìm kiếm cơ bản


%%%%%%%%%%%% Slide %%%%%%%%%%%%%%%%%%%%%%%%%%%%%%%%%%%%%%%%%%%%%%%%%%%
\heading{Tác nhân giải quyết vấn đề}

Một dạng hạn chế của tác nhân tổng quát:

\input{\file{algorithms}{simple-PS-agent-algorithm}}

Lưu ý: đây là giải quyết vấn đề {\em ngoại tuyến (offline)}.\\
Giải quyết vấn đề {\em trực tuyến (online)} bao gồm việc hành động mà không có
tri thức đầy đủ về bài toán và giải pháp.



%%%%%%%%%%%% Slide %%%%%%%%%%%%%%%%%%%%%%%%%%%%%%%%%%%%%%%%%%%%%%%%%%%
\heading{Ví dụ: Romania}

Đang đi nghỉ ở Romania; hiện đang ở Arad.\\
Chuyến bay khởi hành vào ngày mai từ Bucharest

\u{Định dạng mục tiêu}:\nl
có mặt ở Bucharest

\u{Khởi tạo bài toán}:\nl
{\em trạng thái}: các thành phố khác nhau\nl
{\em toán tử}: lái xe giữa các thành phố

\u{Tìm giải pháp}:\nl
chuỗi các thành phố, ví dụ: Arad, Sibiu, Fagaras, Bucharest



%%%%%%%%%%%% Slide %%%%%%%%%%%%%%%%%%%%%%%%%%%%%%%%%%%%%%%%%%%%%%%%%%%
\heading{Ví dụ: Romania}


\epsfxsize=\maxfigwidth
\fig{\file{figures}{romania.ps}}


%%%%%%%%%%%% Slide %%%%%%%%%%%%%%%%%%%%%%%%%%%%%%%%%%%%%%%%%%%%%%%%%%%
\heading{Các loại bài toán}

\u{Tất định, có thể truy cập} $\Longrightarrow$ {\em bài toán trạng thái đơn (single-state)}\\
\u{Tất định, không thể truy cập} $\Longrightarrow$ {\em bài toán đa trạng thái (multiple-state)}

\u{Không tất định, không thể truy cập} $\Longrightarrow$ {\em bài toán dự phòng (contingency)}\nl
phải sử dụng cảm biến trong quá trình thực thi \nl
giải pháp là một {\em cây} hoặc {\em chính sách} \nl
thường {\em đan xen} giữa tìm kiếm và thực thi

\u{Không gian trạng thái chưa biết} $\Longrightarrow$ {\em bài toán khám phá} (``trực tuyến'')



%%%%%%%%%%%% Slide %%%%%%%%%%%%%%%%%%%%%%%%%%%%%%%%%%%%%%%%%%%%%%%%%%%
\heading{Ví dụ: thế giới máy hút bụi}

\begin{tabular}{lr}
\hbox{\begin{minipage}[b]{0.6\textwidth}
\u{Trạng thái đơn}, bắt đầu ở \#5. \q{Giải pháp} 
\\

\u{Đa trạng thái}, bắt đầu ở $\{1,2,3,4,5,6,7,8\}$\\
ví dụ: $Phải$ đi đến $\{2,4,6,8\}$. \q{Giải pháp}
\\

\u{Dự phòng}, bắt đầu ở \#5\\
Định luật Murphy: $Hút$ có thể làm bẩn một tấm thảm đang sạch\\
Cảm biến cục bộ: chỉ biết có bụi và vị trí hiện tại.\\
\q{Giải pháp}
\end{minipage}}
&
\epsfxsize=0.4\textwidth
\epsffile{\file{figures}{vacuum2-space.ps}}
\end{tabular}


%%%%%%%%%%%% Slide %%%%%%%%%%%%%%%%%%%%%%%%%%%%%%%%%%%%%%%%%%%%%%%%%%%
\heading{Khởi tạo bài toán trạng thái đơn}

Một {\em bài toán} được định nghĩa bởi bốn mục:

\u{{\em trạng thái ban đầu}} \quad ví dụ: ``ở Arad''

\u{{\em toán tử}} (hoặc {\em hàm kế tiếp} $S(x)$) \nl
ví dụ: Arad $\rightarrow$ Zerind \qquad Arad $\rightarrow$ Sibiu \qquad v.v.

\u{{\em kiểm tra mục tiêu}}, có thể là\nl
{\em tường minh}, ví dụ: $x$ = ``ở Bucharest''\nl
{\em ngầm định}, ví dụ: $NoDirt(x)$

\u{{\em chi phí đường đi}} (cộng dồn)\nl
ví dụ: tổng khoảng cách, số lượng toán tử đã thực thi, v.v.

Một {\em giải pháp} là một chuỗi các toán tử\\
dẫn từ trạng thái ban đầu đến một trạng thái mục tiêu

%%%%%%%%%%%% Slide %%%%%%%%%%%%%%%%%%%%%%%%%%%%%%%%%%%%%%%%%%%%%%%%%%%
\heading{Chọn lựa một không gian trạng thái}

Thế giới thực thì phức tạp một cách vô lý \nl
$\Rightarrow$ không gian trạng thái phải được {\em trừu tượng hóa} để giải quyết bài toán

(Trừu tượng) trạng thái = tập hợp các trạng thái thực

(Trừu tượng) toán tử = kết hợp phức tạp của các hành động thực\nl
   ví dụ: ``Arad $\rightarrow$ Zerind'' đại diện cho một tập hợp phức tạp\nnl
   các tuyến đường có thể, đường vòng, trạm dừng chân, v.v. \\
Để đảm bảo tính khả thi, \u{bất kỳ} trạng thái thực nào ``ở Arad''\al
đều phải đi được đến {\em một} trạng thái thực nào đó ``ở Zerind''

(Trừu tượng) giải pháp = \nl
   tập hợp các đường đi thực tế là giải pháp trong thế giới thực

Mỗi hành động trừu tượng phải ``dễ dàng'' hơn bài toán ban đầu!




%%%%%%%%%%%% Slide %%%%%%%%%%%%%%%%%%%%%%%%%%%%%%%%%%%%%%%%%%%%%%%%%%%
\heading{Ví dụ: Câu đố 8 ô (8-puzzle)}

\epsfxsize=0.6\textwidth
\fig{\file{figures}{8puzzle.ps}}

\q{trạng thái}\\
\q{toán tử}\\
\q{kiểm tra mục tiêu}\\
\q{chi phí đường đi}




%%%%%%%%%%%% Slide %%%%%%%%%%%%%%%%%%%%%%%%%%%%%%%%%%%%%%%%%%%%%%%%%%%
\heading{Ví dụ: Câu đố 8 ô (8-puzzle)}

\epsfxsize=0.6\textwidth
\fig{\file{figures}{8puzzle.ps}}

\q{trạng thái}: vị trí nguyên của các ô vuông (bỏ qua các vị trí trung gian)\\
\q{toán tử}: di chuyển ô trống sang trái, phải, lên, xuống (bỏ qua kẹt v.v.)\\
\q{kiểm tra mục tiêu}: = trạng thái mục tiêu (được cho trước)\\
\q{chi phí đường đi}: 1 cho mỗi lần di chuyển

[Lưu ý: giải pháp tối ưu của họ Câu đố $n$-ô là NP-khó]



%%%%%%%%%%%% Slide %%%%%%%%%%%%%%%%%%%%%%%%%%%%%%%%%%%%%%%%%%%%%%%%%%%
\heading{Ví dụ: đồ thị không gian trạng thái thế giới máy hút bụi}

\epsfxsize=0.8\textwidth
\fig{\file{figures}{vacuum2-paths.ps}}

\q{trạng thái}\\
\q{toán tử}\\
\q{kiểm tra mục tiêu}\\
\q{chi phí đường đi}


%%%%%%%%%%%% Slide %%%%%%%%%%%%%%%%%%%%%%%%%%%%%%%%%%%%%%%%%%%%%%%%%%%
\heading{Ví dụ: đồ thị không gian trạng thái thế giới máy hút bụi}

\epsfxsize=0.8\textwidth
\fig{\file{figures}{vacuum2-paths.ps}}

\q{trạng thái}: số nguyên biểu thị vị trí bụi và robot (bỏ qua {\em lượng} bụi)\\
\q{toán tử}: $Trái$, $Phải$, $Hút$\\
\q{kiểm tra mục tiêu}: không còn bụi\\
\q{chi phí đường đi}: 1 cho mỗi toán tử


%%%%%%%%%%%% Slide %%%%%%%%%%%%%%%%%%%%%%%%%%%%%%%%%%%%%%%%%%%%%%%%%%%
\heading{Ví dụ: lắp ráp robot}

\epsfxsize=0.5\textwidth
\fig{\file{figures}{stanford-arm+blocks.ps}}

\q{trạng thái}: tọa độ giá trị thực của\nl
các khớp robot\nl
các bộ phận của vật thể cần lắp ráp

\q{toán tử}: chuyển động liên tục của các khớp robot

\q{kiểm tra mục tiêu}: lắp ráp hoàn chỉnh {\em không tính robot!}

\q{chi phí đường đi}: thời gian thực thi


%%%%%%%%%%%% Slide %%%%%%%%%%%%%%%%%%%%%%%%%%%%%%%%%%%%%%%%%%%%%%%%%%%
\heading{Các thuật toán tìm kiếm}

Ý tưởng cơ bản:\al
khám phá ngoại tuyến, mô phỏng không gian trạng thái\al
bằng cách sinh ra các trạng thái kế tiếp từ những trạng thái đã khám phá\nnl
(còn gọi là {\em khai triển} trạng thái)

\input{\file{algorithms}{informal-gs-algorithm}}




%%%%%%%%%%%% Slide %%%%%%%%%%%%%%%%%%%%%%%%%%%%%%%%%%%%%%%%%%%%%%%%%%%
\heading{Ví dụ tìm kiếm tổng quát}

\epsfxsize=0.7\textwidth
\fig{\sfile{figures}{general-romania1.ps}}

%%%%%%%%%%%% Overlay %%%%%%%%%%%%%%%%%%%%%%%%%%%%%%%%%%%%%%%%%%%%%%%%%%%
\pheading{Ví dụ tìm kiếm tổng quát}

\epsfxsize=0.7\textwidth
\fig{\sfile{figures}{general-romania2.ps}}


%%%%%%%%%%%% Overlay %%%%%%%%%%%%%%%%%%%%%%%%%%%%%%%%%%%%%%%%%%%%%%%%%%%
\pheading{Ví dụ tìm kiếm tổng quát}

\epsfxsize=0.7\textwidth
\fig{\sfile{figures}{general-romania3.ps}}


%%%%%%%%%%%% Overlay %%%%%%%%%%%%%%%%%%%%%%%%%%%%%%%%%%%%%%%%%%%%%%%%%%%
\pheading{Ví dụ tìm kiếm tổng quát}

\epsfxsize=0.7\textwidth
\fig{\sfile{figures}{general-romania4.ps}}


%%%%%%%%%%%% Slide %%%%%%%%%%%%%%%%%%%%%%%%%%%%%%%%%%%%%%%%%%%%%%%%%%%
\heading{Cài đặt các thuật toán tìm kiếm}

\input{\file{algorithms}{general-search-algorithm}}


%%%%%%%%%%%% Slide %%%%%%%%%%%%%%%%%%%%%%%%%%%%%%%%%%%%%%%%%%%%%%%%%%%
\heading{Cài đặt (tiếp theo): trạng thái (states) vs.~nút (nodes)}

Một {\em trạng thái} là một (biểu diễn của) cấu hình vật lý\\
Một {\em nút} là một cấu trúc dữ liệu tạo thành một phần của cây tìm kiếm\nl
    bao gồm {\em nút cha}, {\em các nút con}, {\em độ sâu}, {\em chi phí đường đi} $g(x)$\\
{\em Trạng thái} không có nút cha, nút con, độ sâu, hay chi phí đường đi!

\epsfxsize=0.65\textwidth
\fig{\file{figures}{state-vs-node.ps}}

Hàm {\sc Expand} (Khai triển) tạo ra các nút mới, điền vào các trường
khác nhau và sử dụng {\sc Operators} (Toán tử) (hoặc {\sc SuccessorFn}) của
bài toán để tạo ra các trạng thái tương ứng.


%%%%%%%%%%%% Slide %%%%%%%%%%%%%%%%%%%%%%%%%%%%%%%%%%%%%%%%%%%%%%%%%%%
\heading{Các chiến lược tìm kiếm}

Một chiến lược được xác định bằng cách chọn {\em thứ tự khai triển nút}

Các chiến lược được đánh giá dựa trên các khía cạnh sau:\nl
\u{tính hoàn chỉnh (completeness)}---nó có luôn tìm ra giải pháp nếu giải pháp tồn tại không?\nl
\u{độ phức tạp thời gian}---số lượng nút được tạo/khai triển\nl
\u{độ phức tạp không gian}---số lượng nút tối đa lưu trong bộ nhớ\nl
\u{tính tối ưu (optimality)}---nó có luôn tìm ra giải pháp có chi phí thấp nhất không?

Độ phức tạp thời gian và không gian được đo lường thông qua\nl
$b$---hệ số rẽ nhánh (branching factor) tối đa của cây tìm kiếm\nl
$d$---độ sâu của giải pháp có chi phí thấp nhất\nl
$m$---độ sâu tối đa của không gian trạng thái (có thể là $\infty$)


%%%%%%%%%%%% Slide %%%%%%%%%%%%%%%%%%%%%%%%%%%%%%%%%%%%%%%%%%%%%%%%%%%
\heading{Các chiến lược tìm kiếm mù (Uninformed search)}

Chiến lược {\em tìm kiếm mù} chỉ sử dụng những thông tin có sẵn\\
trong định nghĩa của bài toán

Tìm kiếm theo chiều rộng (Breadth-first search)

Tìm kiếm chi phí đồng nhất (Uniform-cost search)

Tìm kiếm theo chiều sâu (Depth-first search)

Tìm kiếm sâu hạn chế (Depth-limited search)

Tìm kiếm sâu lặp sâu dần (Iterative deepening search)


%%%%%%%%%%%% Slide %%%%%%%%%%%%%%%%%%%%%%%%%%%%%%%%%%%%%%%%%%%%%%%%%%%
\heading{Tìm kiếm theo chiều rộng}

Khai triển nút chưa khai triển ở độ sâu nông nhất

\u{Cài đặt}:\nl
{\sc QueueingFn} (Hàm hàng đợi) = đặt các nút kế tiếp vào cuối hàng đợi

\epsfxsize=1.1\textwidth
\fig{\sfile{figures}{bfs-romania1.ps}}


%%%%%%%%%%%% Overlay %%%%%%%%%%%%%%%%%%%%%%%%%%%%%%%%%%%%%%%%%%%%%%%%%%%
\pheading{Tìm kiếm theo chiều rộng}

.

.\\
.

\epsfxsize=1.1\textwidth
\fig{\sfile{figures}{bfs-romania2.ps}}


%%%%%%%%%%%% Overlay %%%%%%%%%%%%%%%%%%%%%%%%%%%%%%%%%%%%%%%%%%%%%%%%%%%
\pheading{Tìm kiếm theo chiều rộng}

.

.\\
.

\epsfxsize=1.1\textwidth
\fig{\sfile{figures}{bfs-romania3.ps}}


%%%%%%%%%%%% Overlay %%%%%%%%%%%%%%%%%%%%%%%%%%%%%%%%%%%%%%%%%%%%%%%%%%%
\pheading{Tìm kiếm theo chiều rộng}

.

.\\
.

\epsfxsize=1.1\textwidth
\fig{\sfile{figures}{bfs-romania4.ps}}


%%%%%%%%%%%% Slide %%%%%%%%%%%%%%%%%%%%%%%%%%%%%%%%%%%%%%%%%%%%%%%%%%%
\heading{Các thuộc tính của tìm kiếm theo chiều rộng}

\q{Hoàn chỉnh}

\q{Thời gian}

\q{Không gian}

\q{Tối ưu}


%%%%%%%%%%%% Slide %%%%%%%%%%%%%%%%%%%%%%%%%%%%%%%%%%%%%%%%%%%%%%%%%%%
\heading{Các thuộc tính của tìm kiếm theo chiều rộng}

\q{Hoàn chỉnh} Có (nếu $b$ là hữu hạn)

\q{Thời gian} $1+b+b^2+b^3+\ldots +b^d = O(b^d)$, tức là tăng theo hàm mũ của $d$

\q{Không gian} $O(b^d)$ (lưu giữ mọi nút trong bộ nhớ)

\q{Tối ưu} Có (nếu chi phí = 1 cho mỗi bước); thông thường thì không tối ưu

{\em Không gian} là một vấn đề lớn; có thể dễ dàng tạo ra các nút ở mức 1MB/giây\nl
do đó 24 giờ = 86GB.


%%%%%%%%%%%% Slide %%%%%%%%%%%%%%%%%%%%%%%%%%%%%%%%%%%%%%%%%%%%%%%%%%%
\heading{Romania với chi phí bước tính bằng km}

\epsfxsize=1.0\textwidth
\fig{\file{figures}{romania2.ps}}


%%%%%%%%%%%% Slide %%%%%%%%%%%%%%%%%%%%%%%%%%%%%%%%%%%%%%%%%%%%%%%%%%%
\heading{Tìm kiếm chi phí đồng nhất}

Khai triển nút chưa khai triển có chi phí thấp nhất

\u{Cài đặt}:\nl
{\sc QueueingFn} = chèn theo thứ tự chi phí đường đi tăng dần

\epsfxsize=1.1\textwidth
\fig{\sfile{figures}{uc-romania1.ps}}



%%%%%%%%%%%% Overlay %%%%%%%%%%%%%%%%%%%%%%%%%%%%%%%%%%%%%%%%%%%%%%%%%%%
\pheading{Tìm kiếm chi phí đồng nhất}

.

.\\
.

\epsfxsize=1.1\textwidth
\fig{\sfile{figures}{uc-romania2.ps}}

%%%%%%%%%%%% Overlay %%%%%%%%%%%%%%%%%%%%%%%%%%%%%%%%%%%%%%%%%%%%%%%%%%%
\pheading{Tìm kiếm chi phí đồng nhất}

.

.\\
.

\epsfxsize=1.1\textwidth
\fig{\sfile{figures}{uc-romania3.ps}}


%%%%%%%%%%%% Overlay %%%%%%%%%%%%%%%%%%%%%%%%%%%%%%%%%%%%%%%%%%%%%%%%%%%
\pheading{Tìm kiếm chi phí đồng nhất}

.

.\\
.

\epsfxsize=1.1\textwidth
\fig{\sfile{figures}{uc-romania4.ps}}



%%%%%%%%%%%% Slide %%%%%%%%%%%%%%%%%%%%%%%%%%%%%%%%%%%%%%%%%%%%%%%%%%%
\heading{Các thuộc tính của tìm kiếm chi phí đồng nhất}

\q{Hoàn chỉnh} Có, nếu chi phí mỗi bước $\geq \epsilon$

\q{Thời gian} Số nút có $g \leq {}$ chi phí của giải pháp tối ưu

\q{Không gian} Số nút có $g \leq {}$ chi phí của giải pháp tối ưu

\q{Tối ưu} Có


%%%%%%%%%%%% Slide %%%%%%%%%%%%%%%%%%%%%%%%%%%%%%%%%%%%%%%%%%%%%%%%%%%
\heading{Tìm kiếm theo chiều sâu}

Khai triển nút chưa khai triển ở độ sâu sâu nhất

\u{Cài đặt}:\nl
{\sc QueueingFn} = chèn các nút kế tiếp vào đầu hàng đợi

\epsfxsize=1.1\textwidth
\fig{\sfile{figures}{dfs-romania1.ps}}


%%%%%%%%%%%% Overlay %%%%%%%%%%%%%%%%%%%%%%%%%%%%%%%%%%%%%%%%%%%%%%%%%%%
\pheading{Tìm kiếm theo chiều sâu}

.

.\\
.

\epsfxsize=1.1\textwidth
\fig{\sfile{figures}{dfs-romania2.ps}}


%%%%%%%%%%%% Overlay %%%%%%%%%%%%%%%%%%%%%%%%%%%%%%%%%%%%%%%%%%%%%%%%%%%
\pheading{Tìm kiếm theo chiều sâu}

.

.\\
.

\epsfxsize=1.1\textwidth
\fig{\sfile{figures}{dfs-romania3.ps}}


%%%%%%%%%%%% Overlay %%%%%%%%%%%%%%%%%%%%%%%%%%%%%%%%%%%%%%%%%%%%%%%%%%%
\pheading{Tìm kiếm theo chiều sâu}

.

.\\
.

\epsfxsize=1.1\textwidth
\fig{\sfile{figures}{dfs-romania4.ps}}

\vspace*{-0.6in}

Tức là, tìm kiếm theo chiều sâu có thể rơi vào các vòng lặp vô hạn\\
Cần một không gian tìm kiếm hữu hạn, không có chu trình (hoặc phải kiểm tra trạng thái lặp lại)


%%%%%%%%%%%% Slide %%%%%%%%%%%%%%%%%%%%%%%%%%%%%%%%%%%%%%%%%%%%%%%%%%%
\heading{Tìm kiếm theo chiều sâu trên cây nhị phân độ sâu 3}

\vspace*{0.3in}

\epsfxsize=1.0\textwidth
\fig{\sfile{figures}{dfs-binary1.ps}}


%%%%%%%%%%%% Overlay %%%%%%%%%%%%%%%%%%%%%%%%%%%%%%%%%%%%%%%%%%%%%%%%%%%
\pheading{Tìm kiếm theo chiều sâu trên cây nhị phân độ sâu 3}

\vspace*{0.3in}

\epsfxsize=1.0\textwidth
\fig{\sfile{figures}{dfs-binary2.ps}}


%%%%%%%%%%%% Overlay %%%%%%%%%%%%%%%%%%%%%%%%%%%%%%%%%%%%%%%%%%%%%%%%%%%
\pheading{Tìm kiếm theo chiều sâu trên cây nhị phân độ sâu 3}

\vspace*{0.3in}

\epsfxsize=1.0\textwidth
\fig{\sfile{figures}{dfs-binary3.ps}}


%%%%%%%%%%%% Overlay %%%%%%%%%%%%%%%%%%%%%%%%%%%%%%%%%%%%%%%%%%%%%%%%%%%
\pheading{Tìm kiếm theo chiều sâu trên cây nhị phân độ sâu 3}

\vspace*{0.3in}

\epsfxsize=1.0\textwidth
\fig{\sfile{figures}{dfs-binary4.ps}}


%%%%%%%%%%%% Overlay %%%%%%%%%%%%%%%%%%%%%%%%%%%%%%%%%%%%%%%%%%%%%%%%%%%
\pheading{Tìm kiếm theo chiều sâu trên cây nhị phân độ sâu 3}

\vspace*{0.3in}

\epsfxsize=1.0\textwidth
\fig{\sfile{figures}{dfs-binary5.ps}}


%%%%%%%%%%%% SLide %%%%%%%%%%%%%%%%%%%%%%%%%%%%%%%%%%%%%%%%%%%%%%%%%%%
\heading{Tìm kiếm theo chiều sâu trên cây nhị phân độ sâu 3 (tiếp)}

\vspace*{0.3in}

\epsfxsize=1.0\textwidth
\fig{\sfile{figures}{dfs-binary6.ps}}


%%%%%%%%%%%% Overlay %%%%%%%%%%%%%%%%%%%%%%%%%%%%%%%%%%%%%%%%%%%%%%%%%%%
\pheading{Tìm kiếm theo chiều sâu trên cây nhị phân độ sâu 3}

\vspace*{0.3in}

\epsfxsize=1.0\textwidth
\fig{\sfile{figures}{dfs-binary7.ps}}


%%%%%%%%%%%% Overlay %%%%%%%%%%%%%%%%%%%%%%%%%%%%%%%%%%%%%%%%%%%%%%%%%%%
\pheading{Tìm kiếm theo chiều sâu trên cây nhị phân độ sâu 3}

\vspace*{0.3in}

\epsfxsize=1.0\textwidth
\fig{\sfile{figures}{dfs-binary8.ps}}


%%%%%%%%%%%% Overlay %%%%%%%%%%%%%%%%%%%%%%%%%%%%%%%%%%%%%%%%%%%%%%%%%%%
\pheading{Tìm kiếm theo chiều sâu trên cây nhị phân độ sâu 3}

\vspace*{0.3in}

\epsfxsize=1.0\textwidth
\fig{\sfile{figures}{dfs-binary9.ps}}


%%%%%%%%%%%% Slide %%%%%%%%%%%%%%%%%%%%%%%%%%%%%%%%%%%%%%%%%%%%%%%%%%%
\heading{Các thuộc tính của tìm kiếm theo chiều sâu}

\q{Hoàn chỉnh}

\q{Thời gian}

\q{Không gian}

\q{Tối ưu}


%%%%%%%%%%%% Slide %%%%%%%%%%%%%%%%%%%%%%%%%%%%%%%%%%%%%%%%%%%%%%%%%%%
\heading{Các thuộc tính của tìm kiếm theo chiều sâu}

\q{Hoàn chỉnh} Không: thất bại trong không gian có độ sâu vô hạn, không gian có chu trình\nl
Sửa đổi để tránh lặp lại các trạng thái trên đường đi\nnl
$\Rightarrow$ hoàn chỉnh trong các không gian hữu hạn

\q{Thời gian} $O(b^m)$: rất tệ nếu $m$ lớn hơn nhiều so với $d$\nl
        nhưng nếu các giải pháp dày đặc, có thể nhanh hơn nhiều so với tìm kiếm theo chiều rộng

\q{Không gian} $O(bm)$, tức là không gian tuyến tính!

\q{Tối ưu} Không


%%%%%%%%%%%% Slide %%%%%%%%%%%%%%%%%%%%%%%%%%%%%%%%%%%%%%%%%%%%%%%%%%%
\heading{Tìm kiếm sâu hạn chế (Depth-limited search)}

= tìm kiếm theo chiều sâu với giới hạn độ sâu là $l$

\u{Cài đặt}:\nl
Các nút ở độ sâu $l$ không có nút con kế tiếp


%%%%%%%%%%%% Slide %%%%%%%%%%%%%%%%%%%%%%%%%%%%%%%%%%%%%%%%%%%%%%%%%%%
\heading{Tìm kiếm sâu lặp sâu dần (Iterative deepening search)}

\input{\file{algorithms}{ids-algorithm}}



%%%%%%%%%%%% Slide %%%%%%%%%%%%%%%%%%%%%%%%%%%%%%%%%%%%%%%%%%%%%%%%%%%
\heading{Tìm kiếm sâu lặp sâu dần $l=0$}

\vspace*{0.3in}

\epsfxsize=1.1\textwidth
\fig{\sfile{figures}{ids-romania1.ps}}

%%%%%%%%%%%% Slide %%%%%%%%%%%%%%%%%%%%%%%%%%%%%%%%%%%%%%%%%%%%%%%%%%%
\heading{Tìm kiếm sâu lặp sâu dần $l=1$}

\vspace*{0.3in}

\epsfxsize=1.1\textwidth
\fig{\sfile{figures}{ids-romania1.ps}}

%%%%%%%%%%%% Overlay %%%%%%%%%%%%%%%%%%%%%%%%%%%%%%%%%%%%%%%%%%%%%%%%%%%
\pheading{Tìm kiếm sâu lặp sâu dần $l=1$}

\vspace*{0.3in}

\epsfxsize=1.1\textwidth
\fig{\sfile{figures}{ids-romania2.ps}}

%%%%%%%%%%%% Slide %%%%%%%%%%%%%%%%%%%%%%%%%%%%%%%%%%%%%%%%%%%%%%%%%%%
\heading{Tìm kiếm sâu lặp sâu dần $l=2$}

\vspace*{0.3in}

\epsfxsize=1.1\textwidth
\fig{\sfile{figures}{ids-romania1.ps}}

%%%%%%%%%%%% Overlay %%%%%%%%%%%%%%%%%%%%%%%%%%%%%%%%%%%%%%%%%%%%%%%%%%%
\pheading{Tìm kiếm sâu lặp sâu dần $l=2$}

\vspace*{0.3in}

\epsfxsize=1.1\textwidth
\fig{\sfile{figures}{ids-romania2.ps}}

%%%%%%%%%%%% Overlay %%%%%%%%%%%%%%%%%%%%%%%%%%%%%%%%%%%%%%%%%%%%%%%%%%%
\pheading{Tìm kiếm sâu lặp sâu dần $l=2$}

\vspace*{0.3in}

\epsfxsize=1.1\textwidth
\fig{\sfile{figures}{ids-romania3.ps}}

%%%%%%%%%%%% Overlay %%%%%%%%%%%%%%%%%%%%%%%%%%%%%%%%%%%%%%%%%%%%%%%%%%%
\pheading{Tìm kiếm sâu lặp sâu dần $l=2$}

\vspace*{0.3in}

\epsfxsize=1.1\textwidth
\fig{\sfile{figures}{ids-romania4.ps}}

%%%%%%%%%%%% Overlay %%%%%%%%%%%%%%%%%%%%%%%%%%%%%%%%%%%%%%%%%%%%%%%%%%%
\pheading{Tìm kiếm sâu lặp sâu dần $l=2$}

\vspace*{0.3in}

\epsfxsize=1.1\textwidth
\fig{\sfile{figures}{ids-romania5.ps}}




%%%%%%%%%%%% Slide %%%%%%%%%%%%%%%%%%%%%%%%%%%%%%%%%%%%%%%%%%%%%%%%%%%
\heading{Các thuộc tính của tìm kiếm sâu lặp sâu dần}

\q{Hoàn chỉnh}

\q{Thời gian}

\q{Không gian}

\q{Tối ưu}


%%%%%%%%%%%% Slide %%%%%%%%%%%%%%%%%%%%%%%%%%%%%%%%%%%%%%%%%%%%%%%%%%%
\heading{Các thuộc tính của tìm kiếm sâu lặp sâu dần}

\q{Hoàn chỉnh} Có

\q{Thời gian} $(d+1)b^0 + d b^1 + (d-1)b^2 + \ldots + b^d = O(b^d)$

\q{Không gian} $O(bd)$

\q{Tối ưu} Có, nếu chi phí bước = 1\nl
Có thể sửa đổi để khám phá trên cây chi phí đồng nhất


%%%%%%%%%%%% Slide %%%%%%%%%%%%%%%%%%%%%%%%%%%%%%%%%%%%%%%%%%%%%%%%%%%
\heading{Tóm tắt}

Khởi tạo bài toán thường yêu cầu trừu tượng hóa các chi tiết của thế giới thực
để định nghĩa một không gian trạng thái có thể khả thi để khám phá

Nhiều chiến lược tìm kiếm mù khác nhau

Tìm kiếm sâu lặp sâu dần chỉ sử dụng không gian tuyến tính\\
và không tốn nhiều thời gian hơn các thuật toán tìm kiếm mù khác



\end{huge}
\end{document}







